\documentclass[../main.tex]{subfiles}

\begin{document}
\section{Examples}
Recall that for each rational prime $p$ we view $\overline{\Q}$ as a subfield of $\overline{\Q_p}$ and denote by $v_p$ the valuation of $\overline{\Q_p}$.
\begin{lemma}
  Let $K$ be a Galois number field and $p$ a rational prime and assume that:
  \begin{enumerate}
    \item $a = \pm 1$.
    \item $b$ is, up to a unit, a non trivial power of $p$.
    \item $p$ does not split in $K$.
    \item \label{item:valuation_coprimality}$
      1 = \gcd\left( 
        k, 
        v_p\left( b \right) \cdot \left[ v_p\left( K \right): \Z \right]
      \right)
    $.
    \item $n > 2$.
  \end{enumerate}
  Then one of two cases occurs
  \begin{enumerate}
    \item $f$ is irreducible over $K$.
    \item $b = \pm 2$ and $f$ factors into $x \pm 1$ times a polynomial that is irreducible over $K$.
  \end{enumerate}
\end{lemma}
\begin{proof}
  Firstly, note that by considering the Newton polygon of $f$ with respect to $p$ it follows that $f$ has $k$ roots (counted with possible multiplicity, here and in the sequel of the proof) whose $p$-adic valuation is $v_p\left( b \right) / k$ and the rest are of $p$-adic valuation $0$.
  Therefore, the $y$-coordinates of the vertices of the Newton polygon of any factor of $f$ in $K\left[ x \right]$ is in $\frac{v_p\left( b \right)}{k} \Z \cap v_p\left( K \right)$ which according to \ref{item:valuation_coprimality} is $v_p\left( b \right)\Z$ and hence any such factor annihilates either non or all of the $k$ positive valuation roots of $f$.
  Let then $f_0 \in K\left[ X \right]$ be an irreducible monic factor of $f$ that annihilates all of the $k$ positive valuation roots of $f$ and let $f_1 = f_0 / f_1$.
  It suffices to show that either $f / f_0$, or $b = \pm 2$ and $f / f_0 = x \pm 1$.
  
  % To do so, we claim, it is enough to prove that all of the roots of $f_1$ are either $1$ or $-1$.
  % Indeed, assume it is proven.
  % If $f_1$ in fact has no roots at all we are done since $f_0$ is irreducible.
  % Otherwise, since $1$ and $-1$ can not both be roots of $f$ and none of them can be a double root of it, it follows that $f_1 = x \pm 1$ and from the triangle inequality it follows that $\left| b \right| \leq 2$ and hence since $b$ was assumed to be a non trivial $p$-power it follows that $b = \pm 2$ as required.

  We will first show that $f / f_0$ is in $\Q\left[ x \right]$ and we will do so by showing that $f_0$ is in $\Q\left[ x \right]$.
  Let $\sigma \in \Gal\left( \overline{\Q}/ \Q \right)$ be an automorphism of $\overline{\Q}$ and denote by $\sigma$ also the induced automorphism of $\overline{\Q}\left[ x \right]$.
  It suffices to show that $\sigma\left( f_0 \right) = f_0$.
  Since $f$ does not split in $K$, $\sigma\left( f_0 \right)$ is also an irreducible factor $f$ with $k$ roots of positive $p$-adic valuation. 
  However, as mentioned $f$ has only $k$ such roots and hence it has at most one such irreducible factor and hence indeed $f_0 = \sigma\left( f_0 \right)$.

  If $f / f_0 = 1$ we are done.
  Assume otherwise and let $f_1$ be some factor of $f / f_0$ that is irreducible in $\Q\left[ x \right]$ and denote by $b_0$ and $b_1$ the free coefficients of $f_0$ and $f_1$ respectively.
  Since $f_0$ annihilates all of the positive $p$-adic valuation roots of $f$ it follows that $v_p\left( b_0 \right) = v_p\left( b \right)$, and hence, since $b_0 b_1$ divides $b$ in $\Z$ and since $b$ is assumed to be a $p$-power up to a unit, it follows that $b_1$ is a rational unit, i.e. that $b_1 = \pm 1$.
  Then $f_1$ has at last one root, say $\alpha$, within the closed complex unit disk.
  By the triangle inequality it thus follows
  \[
    \left| b \right| \leq \left| \alpha^n \right| + \left| a\alpha^k \right| \leq 2
  \]
  and hence since we assume that $b$ is up to a unit a non trivial power of $p$ it follows that $b = \pm 2$, as we set to prove, and hence that 
  \[
    \alpha^k = \alpha^n = \pm 1.
  \]
  Since $1 = \gcd\left( k,n \right)$ it follows that $\alpha = \pm 1$ or equivalently, since $f_1$ is assumed to be irreducible that $f_1 = x \pm 1$.
  Finally, since $1$ and $=-1$ can not both be roots of $f$ and since none of them can be a double root of $f$ it follows that $f / f_0 = x \pm 1$ completing our proof.
\end{proof}

\begin{lemma}
  Let $p$ be a rational prime and assume that 
  \begin{enumerate}
    \item $a = \pm 1$.
    \item $b$ is up to a unit a non trivial power of $p$.
    \item 
  \end{enumerate}
\end{lemma}

\begin{center}
\begin{tabular}{ |c|c|c| } 
  \hline
  $ab$ 
  &
  \makecell{Minimal number fields satisfying \\ the conditions of Lemma \ref{lem:fields_not_embeddable}} 
   & 
   \makecell{
    Next root discriminant realized by a
    \\  
    minimal number field unramified away from $ab$
  } \\
  \hline
  $2$ & $\Q\left( \sqrt{-1} \right)$ & $2^{3/2}$ \\ 
  $3$ & $\Q\left( \sqrt{-3} \right)$ & $3^{4/3}$\\ 
  $4$ & $\Q\left( \sqrt{-1} \right), \Q\left( \sqrt{2} \right), \Q\left( \sqrt{-2} \right)$ & $> 5.989$ \\
  $5$ & $\Q\left( \sqrt{5} \right)$ & $> 5.989$ \\ 
\hline
\end{tabular}
\end{center}
\begin{lemma}
  \label{lem:fields_not_embeddable}
  Assume that $\left| ab \right| \leq 5$ and that $\Q \subsetneq K$ is a number field that is unramified away from $ab$ for which 
  \[
    v_p\left( \rd K \right) \leq v_p\left( ab \right).
  \]
  Then $K$ can not be embedded into $\Q\left[ x \right]/ \left( f \right)$.
\end{lemma}
\begin{proof}
  
\end{proof}

\subsection{$\lcm\left( a,b \right) = 1$}

\subsection{$\lcm\left( a,b \right) = 2$}
\subsubsection{$\left| a \right| = 2, \left| b \right| = 1$}
\begin{lemma}
  Assume that $n > 7$.
  Then one of two cases occurs:
  \begin{enumerate}
    \item $f\left( 1 \right)f\left( -1 \right) \neq 0$ and $f$ is irreducible.
    \item $f\left( 1 \right)f\left( -1 \right) = 0$ and $f$ factors into $x \pm 1$ times an irreducible factor.
  \end{enumerate}
\end{lemma}
\begin{proof}
  From \cite[p. 248]{filaseta2002irreducibility} it follows that the non cyclotomic part of $f$ is irreducible.
  However, if $\zeta$ is a root of unity and 
  $
    \left| \zeta^n \pm 1 \right| = \left| 2\zeta^k \right|
  $
  then 
  $
    \zeta^n,\zeta^k \in \left\{ \pm 1 \right\}
  $
  and hence $\zeta = \pm 1$.
  Moreover, it is easy to see that $\pm 1$ can not both be roots of $f$ and that neither of them is a double root and hence the cyclotomic part of $f$ is either $1$ or $x \pm 1$.
\end{proof}
\begin{lemma}
  $\Q\left( i \right)$ can not be embedded into $\Q\left[ x \right]/ \left( f \right)$.
\end{lemma}
\begin{proof}
  Assume first that $f\left( 1 \right)f\left( -1 \right) \neq 0$ and hence $f$ is irreducible.
  It suffices to show that $f$ has a real root.
  However, for $f$ to not have a real root, $n$ must be even $k$ must be odd and $b$ must be $1$ and it can be verified that in such a case $0 = f\left( 1 \right)f\left( -1 \right)$ which is a contradiction.
  
  Assume now that $f\left( 1 \right)f\left( -1 \right) = 0$ and let $g$ be the non-linear factor of $f$.
  For $g$ to not have a real root, $n$ must be odd, however in this case $f$ is separable over ${\F}_2$ and hence so does $g$ and hence $L$ is an unramified above $2$ and in particular can not contain $\Q\left( i \right)$.
\end{proof}

\subsubsection{$\left| a \right| = 1, \left| b \right| = 2$}
\end{document}